% Part: first-order-logic
% Chapter: completeness
% Section: downward-ls

\documentclass[../../../include/open-logic-section]{subfiles}

\begin{document}

\olfileid[tr]{fol}{com}{dls}
\olsection{Löwenheim--Skolem Teoremi}

Löwenheim--Skolem Teoremi, bir kuramın sonsuz bir modeli varsa en fazla
!!{denumerable} bir modelinin de bulunduğunu söyler. Bu olgunun hemen
çıkan bir sonucu, birinci dereceden mantığın !!a{structure} boyutunun
!!{nonenumerable} olduğunu ifade edememesidir: herhangi bir !!{sentence}
ya da !!{sentence}s kümesi bütün !!{nonenumerable} !!{structure}s içinde
sağlanıyorsa, !!{enumerable} bir yapı içinde de sağlanır.

\begin{thm}
\ollabel{thm:downward-ls} $\Gamma$ tutarlıysa !!a{enumerable} modeli
vardır; yani !!a{structure} içinde sağlanabilirdir ve bu yapının evreni
ya sonlu ya da !!{denumerable} olur.
\end{thm}

\begin{proof}
$\Gamma$ tutarlıysa tamlık teoreminin ispatının verdiği
!!{structure}~$\Struct M$ yapısının evreni~$\Domain{M}$,~$\Lang L$
dilindeki terimler kümesinden daha büyük değildir. Dolayısıyla
$\Struct M$ en fazla !!{denumerable}{}dur.
\end{proof}

\begin{thm}
\ollabel{noidentity-ls} $\Gamma$, özdeşliksiz birinci dereceden mantık
dilindeki tutarlı bir !!{sentence}s kümesiyse !!a{denumerable} modeli
vardır; yani !!a{structure} içinde sağlanabilirdir ve bu yapının evreni
sonsuz ve !!{enumerable} olur.
\end{thm}

\begin{proof}
$\Gamma$ tutarlıysa ve özdeşliğin geçtiği hiçbir tümce içermiyorsa,
tamlık teoreminin ispatının verdiği !!{structure}~$\Struct M$ yapısının
evreni~$\Domain{M}$,~$\Lang L'$ dilindeki terimler kümesinin tam
kendisidir. $\Struct{M}$, $\Trm[L']$ gibi !!{denumerable} niteliktedir.
\end{proof}

\begin{ex}[Skolem Paradoksu]
Zermelo--Fraenkel küme kuramı~$\Log{ZFC}$, kümelerin büyüklükleri
hakkındaki olgular da dahil olmak üzere hemen hemen bütün matematiksel
ifadelerin dile getirilebildiği çok güçlü bir çerçevedir. Örneğin
$\Log{ZFC}$, gerçel sayılar kümesi~$\Real$'in !!{nonenumerable}
olduğunu; herhangi bir kümenin kuvvet kümesinin o kümenin kendisinden
daha büyük olduğunu söyleyen Cantor Teoremini vb. ispatlayabilir.
$\Log{ZFC}$ tutarlıysa bütün modelleri sonsuzdur; dahası hepsi
!!{element}s içerir ve kuram bu öğelerin !!{nonenumerable} olduğunu söyler.
Örneğin doğal sayıların kuvvet kümesinin var olduğuna ilişkin
$\Log{ZFC}$ teoremini doğru kılan eleman böyledir. Löwenheim--Skolem
Teoremi uyarınca $\Log{ZFC}$'nin ayrıca !!{enumerable} modelleri de
vardır: ``!!{nonenumerable}'' kümeler içeren fakat kendileri
!!{enumerable} olan modeller.
\end{ex}

\end{document}
